%%%%%%%%%%%%%%%%%%%%%%%%%%%%%%%%%%%%%%%%%%%%%%%%%%%%%%%%%%%%%%%%%%%%%%
% Lớp 11 | Dãy số. Cấp số cộng. Cấp số nhân | Cấp số cộng
% Dạng 1D2?2-5: Điều kiện để dãy số là cấp số cộng
% Số câu: 45  (sinh tự động bởi tach_cau.py; từ nay thầy sửa trực tiếp file này)
%%%%%%%%%%%%%%%%%%%%%%%%%%%%%%%%%%%%%%%%%%%%%%%%%%%%%%%%%%%%%%%%%%%%%%

%%%%%%%%%%%%%%%% TRẮC NGHIỆM 4 PHƯƠNG ÁN %%%%%%%%%%%%%%%%

% nguồn: Ok_Toan11_TracNghiem_Lan2
\begin{ex}%[1D2H2-5]%[ma:fb018ba3]
	 Viết ba số xen giữa $2$ và $22$ để ta được một cấp số cộng có $5$ số hạng?
	 \choice
		{$6$, $12$, $18$} 
		{$8$, $13$,  $18$}
		{\True $7$, $12$, $17$} 
		{$6$, $10$,  $14$}
	 \loigiai{
		Xem cấp  số cộng cần tìm là $(u_n)$ có
		 $\heva{&u_1=2\\&u_5=22}$.\\
		 Suy ra $d=5$.\\
		Vậy cấp số cộng cần tìm là $(u_n)$:  $2$, $7$, $12$, $17$, $22$.
	 }
\end{ex}

% nguồn: Ok_Toan11_TracNghiem_Lan2
\begin{ex}%[1D2H2-5]%[ma:02013f3d]
	Xác định $a$ để $3$ số $1+2a$; $2a^2-1$; $-2a$ theo thứ tự thành lập một cấp số cộng? 
	\choice 
		{Không có giá trị nào của  $a$}
		{$a=\pm\dfrac{\sqrt  3}{4}$}
		{$a=\pm  3$}
		{\True  $a=\pm\dfrac{\sqrt 3}{2}$}
	\loigiai{ 
	Theo công thức cấp số  cộng ta có \[2(2a^2-1)=(1+2a)+(-2a)\Leftrightarrow{a^2}=\dfrac{3}{4}\Leftrightarrow a=\pm\dfrac{\sqrt 3}{2}.\]}
\end{ex}

% nguồn: Ok_TracNghiemLop11_Lan1
\begin{ex}%[1D2H2-5]%[ma:8cc3bdd0]
	Cho phương trình $x^3-3x^2-9x+m=0$. Tìm $m$ để phương trình có $3$ nghiệm lập thành cấp  số cộng.
	 \choice
		 {$16$}
		 {$13$}
		{$12$} 
		{\True  11}
	 \loigiai{
		Giả sử phương trình có $3$ nghiệm lập  thành cấp số cộng.\\
		Khi  đó
		\[x_1+x_3=2x_2; x_1+x_2+x_3=3\Rightarrow x_2=1\Rightarrow m=11. 
		\] 
		Ta  có phương trình $x^3-3x^2-9x+11=0\Rightarrow x_1=1-\sqrt{12}$; $x_2=1$; $x_3=1+\sqrt{12}$.\\
		Ba nghiệm này lập thành cấp số  cộng.\\
		 Vậy $m=11$.
	} 
\end{ex}

% nguồn: Ok_TracNghiemLop11_Lan1
\begin{ex}%[1D2H2-5]%[ma:cbe4d4ae]
	Gọi $S$ là tập hợp tất cả các số tự nhiên $k$ sao cho $\mathrm{C}_{14}^k$, $\mathrm{C}_{14}^{k+1}$, $\mathrm{C}_{14}^{k+2}$  theo thứ tự đó lập thành một cấp số cộng. Tính tổng tất cả các phần tử của $S$.
	\choice 
		 {$8$}
		{$6$} 
		 {$10$}
		 {\True $12$}
	 \loigiai{
		 Ta có
		\begin{eqnarray*}
			\mathrm{C}_{14}^k+\mathrm{C}_{14}^{k+2}=2\mathrm{C}_{14}^{k+1}&\Leftrightarrow & f(x)=g(x)\\ 
			&\Leftrightarrow &  \dfrac{14!}{k!\left(14-k\right)!}+\dfrac{14!}{\left(k+2\right)!\left(12-k\right)!}=2\dfrac{14!}{\left(k+1\right)!\left(13-k\right)!}\\
			&\Leftrightarrow  & \dfrac{1}{\left(14-k\right)\left(13-k\right)}+\dfrac{1}{\left(k+1\right)\left(k+2\right)}=\dfrac{2}{\left(k+1\right)\left(13-k\right)}\\
			 &\Leftrightarrow& \left(k+1\right)\left(k+2\right)+\left(14-k\right)\left(13-k\right)=2\left(k+2\right)\left(14-k\right)\\
			&\Leftrightarrow& k^2-12k+32=0  \Leftrightarrow \hoac{&{k=8} \\
			&{k=4.}} 
		\end{eqnarray*}
		Vậy $S=\{4;8\}$. Khi đó tổng các giá  trị trong $S$ là $4+8=12$.
	 }
\end{ex}

% nguồn: Ok_TracNghiemLop11_Lan1
\begin{ex}%[1D2H2-5]%[ma:fb2cdac6]
	 Bốn số $x$, $-2$, $y$, $6 $ theo thứ tự đó lập thành một cấp số cộng. Khẳng định nào sau đây đúng?
	\choice 
		{$x=-6 $;  $y=3 $}
		{$ x= -5$; $y=  3$}
		{\True $x=-6  $; $y=2 $}
		{$ x= -5$; $y=  2$}
	\loigiai{ 
		 Ta có bốn số $x,-2, y, 6$ theo thứ tự lập thành một cấp số cộng khi và chỉ khi $$\heva{&{x+y=-4} \\
		&{2y=-2+6}} \Leftrightarrow  \heva{&{x=-6} \\
		 &{y=2}.}$$
	 }
\end{ex}

% nguồn: Ok_TracNghiemLop11_Lan1
\begin{ex}%[1D2H2-5]%[ma:d2247509]
	Có tất cả bao nhiêu bộ số nguyên dương $\left(n, k\right)$ biết $n < 20$ và các  số $C_n^{k-1}$, $C_n^k$, $C_n^{k+1}$ theo thứ tự đó là số hạng thứ nhất, thứ ba, thứ năm của một cấp số cộng.
	\choice 
		{\True  $4$}
		 {$2$}
		 {$1$}
		{$z_1=a+5i$} 
	\loigiai{ 
		Các số $C_n^{k-1}$, $z_2=b$, $C_n^{k+1}$ theo thứ tự đó là số hạng thứ nhất, thứ ba, thứ  năm của một cấp số cộng nên ta có
		$$C_n^k-C_n^{k-1}=C_n^{k+1}-C_n^k\Leftrightarrow  \dfrac{n!}{\left(k+1\right)!(n-k-1)!}+\dfrac{n!}{\left(k-1\right)!(n-k+1)!}=2\dfrac{n!}{k!(n-k)!}$$
		\[\Leftrightarrow \dfrac{1}{k\left(k+1\right)}+\dfrac{1}{\left(n-k+1\right)\left(n-k\right)}=\dfrac{2}{k\left(n-k\right)} \Leftrightarrow  \left(n-2k\right)^2=n+2.
		\] 
		Do $n  < 20\Rightarrow n+2< 22$ mà $n+2$ là số chính phương, $n$, $k$ nguyên dương nên có các trường hợp sau
		\begin{enumerate}[TH 1.]
			 \item $n+2=4\Rightarrow n=2$; $k=2$.
			\item  $n+2=9\Rightarrow n=7$; $k=2$ hoặc $n=7$; $k=5$.
			\item $n+2=16\Rightarrow  n=14$; $k=5$ hoặc $n=14$; $k=9$.
		\end{enumerate}
		Mà  $k+1\le n$ nên chỉ có $4$ bộ thỏa mãn.
	} 
\end{ex}

% nguồn: Ok_TracNghiemLop11_Lan1
\begin{ex}%[1D2H2-5]%[ma:2ac345fb]
	Cho $a$, $b$, $c$ theo thứ tự  này là ba số hạng liên tiếp của một cấp số cộng. Biết $a+b+c=15$. Giá trị của $b$ bằng
	\choice 
		 {$10$}
		 {$8$}
		 {\True $5$}
		{$6$} 
	 \loigiai{
		$a$, $b$, $c$ theo thứ tự là cấp số cộng và  $a+b+c=15\Rightarrow 3b=15\Rightarrow b=5$.
	 }
\end{ex}

% nguồn: Ok_TracNghiemLop11_Lan1
\begin{ex}%[1D2H2-5]%[ma:09db0ad4]
	Biết các số  $C_n^1$; $C_n^2$; $C_n^3$ theo thứ tự lập thành một cấp số cộng với $n > 3$. Tìm $n.$
	 \choice
		 {$n=5$}
		{\True $n=7$} 
		{$n=9$} 
		{$n=11$} 
	\loigiai{ 
		Ba số $C_n^1$; $C_n^2$; $C_n^3$ theo thứ tự  $u_1$, $u_2$, $u_3$ lập thành cấp số cộng nên
		\[u_1+u_3=2u_2 \Leftrightarrow C_n^1+C_n^3=2C_n^2 \left(n\ge 3\right)\Leftrightarrow n+\dfrac{\left(n-2\right)\left(n-1\right)n}{6}=2\cdot\dfrac{\left(n-1\right)n}{2} 
		 \]
		\[\Leftrightarrow 1+\dfrac{n^2-3n+2}{6}=n-1\Leftrightarrow n^2-9n+14\Leftrightarrow \hoac{&{n=2} \\ 
		&{n=7}} \Leftrightarrow n=7 \left(n\ge 3\right). 
		\] 
		\textbf{Nhận xét} Nếu $u_{k-1}$, $u_k$, $u_{k+1}$ là ba số hạng liên  tiếp của một cấp số cộng thì ta có $u_{k-1}+u_{k+1}=2u_k$.
	} 
\end{ex}

% nguồn: Ok_TracNghiemLop11_Lan1
\begin{ex}%[1D2H2-5]%[ma:8ba77e49]
	Tìm tất cả các số thực $x$ để ba số $x^2$, $x^2+1$, $3x$ theo thứ tự đó lập thành  cấp số cộng?
	\choice 
		 {$x=2$}
		{\True $x\in  \left\{1;2\right\}$}
		{$x=0$} 
		{$x\in  \left\{2;3\right\}$}
	 \loigiai{
		Do ba số $x^2$, $x^2+1$, $3x$ lập  thành cấp số cộng nên $$2\left(x^2+1\right)=x^2+3x\Leftrightarrow x^2-3x+2=0\Leftrightarrow \hoac{&{x=1} \\
		&{x=2}.}$$ 
		 Vậy $x\in \left\{1;2\right\}$.
	} 
\end{ex}

% nguồn: Ok_TracNghiemLop11_Lan1
\begin{ex}%[1D2H2-5]%[ma:09078c25]
	Tìm giá trị của $x$, $y$ sao cho dãy số $-2$, $x$, $4$, $y$ theo thứ  tự lập thành một cấp số cộng.
	 \choice
		{$x=2$, $y=8$} 
		{\True $x=1$,  $y=7$}
		 {$x=2$, $y=10$}
		{$x=-6$,  $y=2$}
	\loigiai{ 
		Dãy số $-2$, $ x$, $ 4$, $ y$ theo thứ tự lập thành  một cấp số cộng $\Rightarrow \heva{x&=\dfrac{-2+4}{2}\\4&=\dfrac{x+y}{2}}\Leftrightarrow \heva{x&=1\\y&=7}$.
	} 
\end{ex}

% nguồn: Ok_TracNghiemLop11_Lan1
\begin{ex}%[1D2H2-5]%[ma:9ce3adb6]
	Cho cấp số cộng $\left(u_n\right)\colon 3$, $a$, $13$, $b$.  Tích $ab$ bằng
	\choice 
		 {\True $144$}
		 {$39$}
		 {$26$}
		 {$104$}
	 \loigiai{
		Ta có $\heva{& 3+13=2a\\& a+b=2.13 
		 }\Leftrightarrow\heva{& a=8\\&b=18
		}\Rightarrow ab=144$. 
	 }
\end{ex}

% nguồn: Ok_TracNghiemLop11_Lan1
\begin{ex}%[1D2H2-5]%[ma:d38db667]
	Phương trình $x^3+ax+b=0$ có ba nghiệm lập thành cấp số cộng  khi và chỉ khi
	 \choice
		{\True $b=0$,  $a < 0$}
		{$b=0$,  $a=1$}
		 {$b=1$, $a=-2$}
		{$b=-2$, $a=1$} 
	 \loigiai{
		Gỉa sử phương  trình có ba nghiệm lập thành cấp số cộng khi đó ta có
		\[x_1+x_3=2x_2, x_1+x_2+x_3=0 \Leftrightarrow x_2=0. 
		\] 
		Thay $x_2=0$ vào phương trình $b=0$. Ta có $b=0$ $\Rightarrow x^3+ax=0$ $\Rightarrow a < 0$. $\Rightarrow  b=0$, $a < 0$ thỏa mãn.
	 }
\end{ex}

% nguồn: Ok_Toan11_TracNghiem_Lan2
\begin{ex}%[1D2V2-5]%[ma:e2985c22]
	Cho cấp số cộng $u_1$; $u_2$; $u_3$; $\dots$;  $u_n$ có công sai $d$, các số hạng của cấp số cộng đã cho đều khác $0$.
	Với giá trị nào của  $d$ thì dãy số
	$\dfrac{1}{u_1}$;  $\dfrac{1}{u_2}$; $\dfrac{1}{u_3}$; $\dots$; $\dfrac{1}{u_n}$
	 là một cấp số cộng?
	\choice 
		{$d=-1$} 
		{\True $d=0$} 
		 {$d=1$}
		{$d=2$} 
	 \loigiai{
		Ta có 
		 \[
		 \heva{
		&u_2  - u_1 = d\\
		&u_3 - u_2 =  d
		 }
		\Rightarrow 
		\heva{ 
		&\dfrac{1}{u_2} - \dfrac{1}{u_1} = -\dfrac{d}{u_1u_2}  \\
		 &\dfrac{1}{u_3} - \dfrac{1}{u_2} = -\dfrac{d}{u_2u_3}
		 }.
		\] 
		Vì dãy số $\dfrac{1}{u_1}$; $\dfrac{1}{u_2}$; $\dfrac{1}{u_3}$; $\dots$; $\dfrac{1}{u_n}$ là một cấp số cộng  nên ta có
		\begin{eqnarray*}
			&&\dfrac{1}{u_2} -  \dfrac{1}{u_1} = \dfrac{1}{u_3} - \dfrac{1}{u_2}\\
			 &\Leftrightarrow &
			\hoac{ 
			&d =  0\\
			 &\dfrac{1}{u_1} = \dfrac{1}{u_3}
			}\\ 
			 &\Leftrightarrow&
			\hoac{ 
			&d  = 0\\
			&u_1 =  u_3 = u_1 + 2d
			 }\\
			 &\Leftrightarrow& d = 0.
		\end{eqnarray*}
	 }
\end{ex}

% nguồn: Ok_TracNghiemLop11_Lan1
\begin{ex}%[1D2V2-5]%[ma:aaf58787]
	Cho cấp số cộng $(u_n)$ có $u_{2013}+u_6=1000$. Tổng $2018$ số hạng đầu tiên của cấp số cộng đó là 
	 \choice
		{\True  $1\,009\,000$}
		 {$100\,800$}
		 {$1\,008\,000$}
		 {$100\,900$}
	\loigiai{ 
		Gọi $d$ là công sai của cấp số cộng.\\ 
		Khi  đó
		\begin{eqnarray*}
			&  & u_{2013}+u_6=1\,000\\
			 &\Leftrightarrow & u_1+2012d+u_1+5d=1\,000\\
			&\Leftrightarrow & 2u_1+2017\cdot  d=1\,000
		\end{eqnarray*}
	Ta có $S_{2018}=2\,018\cdot u_1+\dfrac{2\,017\cdot 2\,018}{2}\cdot  d=1\,009\cdot \left(2u_1+2017\cdot d\right)=1\,009\,000$.}
\end{ex}


%%%%%%%%%%%%%%%% ĐÚNG/SAI %%%%%%%%%%%%%%%%

% nguồn: Ok_dungsai_lop11_Lan2
\begin{ex}%[1D2N2-5]%[ma:cb6c22b5]
	Xác định tính đúng, sai  của các khẳng định sau?
	 \choiceTF
		{\True Dãy  số $(u_n)$ với $\dfrac{-2}{3};\,\dfrac{-1}{3};\, 0;\,\dfrac{1}{3};\,\dfrac{2}{3};\, 1;\,\dfrac{4}{3}$ là cấp số cộng với ${u_1=\dfrac{-2}{3};\, d=\dfrac{1}{3}}$}
		{\True Dãy  số $(u_n)$ với ${u_n=7-3n}$ là cấp số cộng với $u_1=4;\, d=-3$}
		{Dãy số $(u_n)$ với ${u_n=n^2+n+1}$ là cấp số cộng với $u_1=3;\, d=1$} 
		{\True Dãy số $(u_n)$ với $u_n={(-1)^n}+3n$ không là cấp  số cộng}
	 \loigiai{
		\begin{itemchoice}
			\itemch Xét dãy số $(u_n)$ với ${\dfrac{-2}{3};\,\dfrac{-1}{3};  0;\,\dfrac{1}{3};\,\dfrac{2}{3}; 1;\,\dfrac{4}{3}}$.\\
			Ta thấy ${u_2-u_1=u_3-u_2=u_4-u_3=\ldots=\dfrac{1}{3}}$.\\ Vậy $(u_n)$ là cấp số cộng với ${u_1=\dfrac{-2}{3};\,  d=\dfrac{1}{3}}$.
			\itemch Xét dãy  số $(u_n)$ với ${u_n=7-3n}$.\\
			Ta có ${u_{n+1}}-u_n=\left[7-3(n+1)\right]-(7-3n)=-3$.\\ 
			Vậy $(u_n)$ là cấp số cộng  với ${u_1=7-3\cdot 1=4;\, d=-3}$.
			\itemch Xét  dãy số $(u_n)$ với ${u_n=n^2+n+1}$. \\
			Ta có ${u_{n+1}}-u_n={(n+1)^2}+(n+1)+1-\left(n^2+n+1\right)=2n+2$  phụ thuộc vào $n$. \\
			Vậy $(u_n)$ không là cấp  số cộng.
			\itemch Xét dãy  số $(u_n)$ với ${u_n=(-1)^n+3n}$. \\
			 Ta có
			\begin{align*}
				 u_{n+1}-u_n&=(-1)^{n+1}+3(n+1)-\left[{{(-1)}^n}+3n \right]\\
				&=-{(-1)^n}+3-{(-1)^n}\\ 
				 &=3-2(-1)^n.
			\end{align*}
			phụ thuộc vào $n$.\\  Vậy $(u_n)$ không là cấp số cộng.
		\end{itemchoice}
	 }
\end{ex}

% nguồn: Ok_dungsai_lop11_Lan2
\begin{ex}%[1D2N2-5]%[ma:dd7fa946]
	Cho các dãy số $(a_n)$; $(b_n)$; $(c_n)$ có số hạng tổng quát lần  lượt là $a_n=4n-3$; $b_n=\dfrac{2-3n}{4}$; $c_n=n^2$.
	 \choiceTF
		{\True $(a_n)$ là một cấp số cộng với số  hạng đầu $a_1=1$}
		{\True $(a_n)$ là một cấp số cộng với công sai $d=4$} 
		{$(b_n)$ là một cấp số cộng với số hạng đầu $b_1=-\dfrac{1}{4}$ và công sai $d=\dfrac{3}{4}$} 
		 {$(c_n)$ là một cấp số cộng với công sai $d=2$}
	\loigiai{ 
		\begin{itemchoice}
			 \itemch $a_1={4\cdot 1-3}=1$.
			\itemch Ta có $a_{n+1}-a_n=4(n+1)-3-(4n-3)=4,\forall n\ge 1$.  \\
			Do đó $(a_n)$ là một cấp số cộng với số hạng đầu $a_1=4\cdot  1-3=1$ và công sai ${d=4}$.
			\itemch Ta có $b_{n+1}-b_n=\dfrac{2-3(n+1)}{4}-\dfrac{2-3  n}{4}=\dfrac{2-3 n-3-2+3n}{4}=-\dfrac{3}{4},\forall n\geq 1$. \\
			Suy ra$(b_n)$ là một cấp số  cộng với số hạng đầu $b_1=\dfrac{2-3\cdot1}{4}=-\dfrac{1}{4}$ và công sai $d=-\dfrac{3}{4}$.\\
			\itemch  Ta có $c_{n+1}-c_n=(n+1)^2-n^2=2 n+1$ (phụ thuộc vào giá trị của $n$). \\
			Suy ra $(c_n)$  không phải là một cấp số cộng.
		\end{itemchoice}
	} 
\end{ex}

% nguồn: Ok_dungsai_lop11_Lan1
\begin{ex}%[1D2H2-5]%[ma:e3906e94]
 Cho cấp số cộng $(u_n)$  biết $\heva{&u_{7} - u_{1} = 12\\&3u_{3} - 2u_{5} = 7}$. Các mệnh đề sau đúng hay sai?
  \choiceTF
  {\True Cấp số cộng có số hạng đầu là một số  nguyên tố}
   {\True Cấp số cộng có công sai $d=2$}
  {Số  hạng $u_6=20$}
  {Tổng $20$ số hạng đầu của cấp số cộng bằng $500$} 
  \loigiai{
  \begin{itemchoice}
    \itemch Ta có
   $$\heva{&u_{7} - u_{1} = 12\\&3u_{3} - 2u_{5} = 7}\Leftrightarrow \heva{&u_{1} + 6d - u_{1} = 12\\&3u_{1} +  6d - 2u_{1} - 8d = 7}\Leftrightarrow \heva{&d = 2\\&u_{1} - 2d = 7}\Leftrightarrow \heva{&d = 2\\&u_{1} = 11.}$$
    Vậy số hạng đầu $u_1=11$ và công sai $d=2$.
   \itemch Do đó $u_{6}=u_1+5d=11+5\cdot2=21$. 
   \itemch Ta có $S_{n}=nu_1+\dfrac{n(n-1)d}{2}$.  Suy ra $S_{20}=20\cdot11+\dfrac{20(20-1)\cdot2}{2}=600$.
  \end{itemchoice}
  }
\end{ex}

% nguồn: Ok_dungsai_lop11_Lan1
\begin{ex}%[1D2H2-5]%[ma:1fc49abb]
 Cho cấp số cộng $(u_n)$ biết $\heva{& u_2 - u_3 + u_5 =  10 \\ &u_1 + u_6 = 17}$. Các mệnh đề sau đúng hay sai?
  \choiceTF
   {Cấp số cộng có $u_1=-1$, công sai $d=3$}
  {Số hạng  $u_4$ bằng $8$}
   {\True Số hạng tổng quát $u_n=3n-2$}
  {\True Tổng $15$ số hạng đầu  của cấp số cộng bằng $330$}
  \loigiai{
  \begin{itemchoice}
   \itemch Ta có $\heva{&u_2 - u_3 + u_5 = 10 \\ &u_1 + u_6 = 17} \Leftrightarrow \heva{&u_1 + d - u_1 - 2d + u_1 + 4d = 10 \\ &u_1 + u_1 + 5d = 17} \Leftrightarrow \heva{&u_1 + 3d = 10 \\ &2u_1 + 5d = 17} \Leftrightarrow \heva{&u_1 = 1 \\ &d =  3.}$\\
   Vậy $u_1  = 1$ và $d = 3$.
   \itemch  Ta có $u_4=u_1+3d=1+3\cdot 3=10$.
   \itemch Số hạng tổng quát $u_n=u_1+(n-1)d=1+(n-1)\cdot  3=3n-2$.
   \itemch Ta có  $S_{n}=nu_1+\dfrac{n(n-1)d}{2}$. Suy ra $S_{15}=15\cdot1+\dfrac{15(15-1)\cdot3}{2}=330$.
  \end{itemchoice}
  }
\end{ex}

% nguồn: Ok_dungsai_lop11_Lan1
\begin{ex}%[1D2H2-5]%[ma:a1ad63f8]
 Cho cấp số cộng $\left(u_n\right)$ có $u_1=2$,  $d=4$. Các mệnh đề sau đúng hay sai?
  \choiceTF
  {\True Số hạng $u_2$ bằng  $6$}
   {\True Số hạng $u_7$ bằng $26$}
  {Số hạng tổng quát  $u_n=2n-4$}
   {\True Tổng $S=u_7+u_9+u_{11}+\cdots+u_{2017}=4070276$}
  \loigiai{
  \begin{itemchoice}
   \itemch Ta có $u_2=u_1+d=2+4=6$. 
   \itemch Ta có $u_7=u_1+6d=2+6\cdot 4=26$. 
   \itemch Ta có $u_n=u_1+(n-1)d=2+(n-1)\cdot4=4n-2$. 
   \itemch Các số $u_7, u_9, u_{11},\ldots,u_{2017}$ lập thành cấp số cộng có $1006$ số hạng với số hạng đầu $u_7$ công sai là $2d$. 
   $$S=S_{1006}=\dfrac{1006}{2}\left[2u_7+(1006-1)2d\right]=503(2\cdot 26+1005\cdot 8)=4070276.$$ 
  \end{itemchoice}
 } 
\end{ex}

% nguồn: Ok_dungsai_lop11_Lan1
\begin{ex}%[1D2H2-5]%[ma:376ee526]
 Cho cấp số cộng $(u_n)$ biết $\heva{&u_3+u_7=22\\&u_9-u_2=28}$. Các mệnh đề sau đúng hay  sai?
 \choiceTF 
  {\True Cấp số cộng có số hạng đầu bằng $-5$} 
  {Cấp số cộng có  công sai $d=-4$}
  {\True Số hạng $u_5=11$} 
  {Tổng $15$ số hạng  đầu của cấp số cộng bằng $350$}
 \loigiai{ 
  \begin{itemchoice}
   \itemch Ta  có
    $$\heva{&u_3+u_7=22\\&u_9-u_2=28}\Leftrightarrow \heva{&u_1+2d+u_1+6d=22\\&u_1+8d-(u_1+d)=28} \Leftrightarrow \heva{&2u_1+8d=22\\&7d=28}\Leftrightarrow\heva{&u_1=-5\\&d=4.}$$
   Vậy số hạng đầu $u_1=-5$ và công  sai $d=4$.
   \itemch Do đó  $u_{5}=u_1+4d=-5+4\cdot4=11$.
   \itemch Ta có $S_{n}=nu_1+\dfrac{n(n-1)d}{2}$. Suy  ra $S_{15}=15\cdot(-5)+\dfrac{15(15-1)\cdot4}{2}=345$.
  \end{itemchoice}
 } 
\end{ex}

% nguồn: Ok_dungsai_lop11_Lan1
\begin{ex}%[1D2H2-5]%[ma:bc6cf20b]
 Cho cấp số cộng $(u_n)$ có $u_1=-2$, $u_6=8$. Các mệnh đề sau  đúng hay sai?
  \choiceTF
  {Công sai $d$ bằng  $3$}
  {Số hạng $u_9$ bằng  $12$}
  {\True Số hạng tổng quát $u_n=2n-4$} 
  {\True Tổng $40$ số hạng đầu  của cấp số cộng bằng $1480$}
 \loigiai{ 
  \begin{itemchoice}
   \itemch Ta có  $u_6=u_1+5d\Leftrightarrow 8=-2+5\cdot d\Leftrightarrow d=2$.
   \itemch Ta có $u_9=u_1+8d=-2+8\cdot  2=14$.
   \itemch Số hạng tổng quát $u_n=u_1+(n-1)d=-2+(n-1)\cdot2=2n-4$. 
   \itemch Ta có  $S_{n}=nu_1+\dfrac{n(n-1)d}{2}$. Suy ra $S_{40}=40\cdot(-2)+\dfrac{40(40-1)\cdot2}{2}=1480$.
  \end{itemchoice}
 } 
\end{ex}

% nguồn: Ok_dungsai_lop11_Lan1
\begin{ex}%[1D2H2-5]%[ma:5e7d2007]
 Cho cấp số cộng $(u_n)$ biết $\heva{&u_1+u_6=10\\&S_{15}=165}$. Các mệnh  đề sau đúng hay sai?
  \choiceTF
  {Cấp số cộng có số hạng  đầu là số nguyên}
  {\True Cấp số cộng  có công sai là số hữu tỉ}
  {\True Số hạng  $u_{10}=\dfrac{44}{3}$}
  {\True Tổng $20$ số hạng đầu của  cấp số cộng bằng $\dfrac{860}{3}$}
 \loigiai{ 
  \begin{itemchoice}
   \itemch Ta có $\heva{&u_1+u_6=10\\&S_{15}=165}\Leftrightarrow \heva{&2u_1+5d=10\\&2u_1+14d=22}\Leftrightarrow  \heva{&u_1=\dfrac{5}{3}\\&d=\dfrac{4}{3}.}$
   \itemch Do đó  $u_{10}=u_1+9d= \dfrac{5}{3}+9\cdot \dfrac{4}{3}=\dfrac{44}{3}$.
   \itemch Ta có $S_{20}=\dfrac{20}{2}\left( 2u_1+19d\right)  = 10\left( \dfrac{10}{3}+\dfrac{76}{3}\right) = \dfrac{860}{3}$.
  \end{itemchoice}
  }
\end{ex}

% nguồn: Ok_dungsai_lop11_Lan1
\begin{ex}%[1D2H2-5]%[ma:d19957de]
  Cho cấp số cộng $(u_n)$ có $u_1=1$, công sai $d=-2$. Các mệnh đề sau đúng hay sai?
 \choiceTF 
  {Số hạng thứ  ba bằng $-2$}
  {Số $-11$ là số hạng thứ  $6$}
  {Số hạng tổng quát là  $u_n=1-2n$}
  {\True Tổng $10$ số hạng đầu bằng $-80$} 
  \loigiai{
  \begin{itemchoice}
   \itemch Ta có  $u_3=u_1+2d=1+2\cdot(-2)=-3$.
   \itemch Giả sử $-11$ là  số hạng thứ $n$ của cấp số cộng.\\
   Ta có $n=\dfrac{u_n-u_1}{d}+1=\dfrac{-11-1}{-2}+1=7$.\\ 
   Vậy số $-11$ là số hạng thứ  $7$ của cấp số cộng $(u_n)$.
    \itemch Số hạng tổng quát $u_n=u_1+(n-1)d=1+(n-1)\cdot(-2)=3-2n$.
   \itemch Ta có $S_{n}=nu_1+\dfrac{n(n-1)d}{2}$.  Suy ra $S_{10}=10\cdot1+\dfrac{10(10-1)\cdot(-2)}{2}=-80$.
  \end{itemchoice}
  }
\end{ex}

% nguồn: Ok_dungsai_lop11_Lan1
\begin{ex}%[1D2H2-5]%[ma:7a072011]
 Cho cấp số cộng $\left(u_n\right)$, biết $\heva{& u_2-3u_5+2u_7=3\\& u_1+2u_4-2u_6=-7}$. Các mệnh đề sau đúng hay  sai?
  \choiceTF
  {\True Cấp  số cộng có $u_1=5$}
  {\True Cấp số cộng có công sai  $d=-3$}
  {Số hạng tổng quát của $\left(u_n\right)$ là $u_n=3n-2$} 
  {\True Tổng của $20$ số hạng đầu tiên của  cấp số cộng bằng $670$}
 \loigiai{ 
  \begin{itemchoice}
   \itemch  Ta có
   $$\heva{& u_2-3u_5+2u_7=3\\& u_1+2u_4-2u_6=-7}\Leftrightarrow  \heva{& (u_1+d)-3(u_1+4d)+2(u_1+6d)=3\\& u_1+2(u_1+3d)-2(u_1+5d)=-7}\Leftrightarrow \heva{& d=3\\& u_1=5.}$$
   \itemch  Ta có $u_n=u_1+(n-1)d=5+(n-1)\cdot3=3n+2$.
   \itemch Ta  có $S_{20}=\dfrac{20}{2}(2u_1+19d)=670$.
  \end{itemchoice}
  }
\end{ex}

% nguồn: Ok_dungsai_lop11_Lan1
\begin{ex}%[1D2H2-5]%[ma:baea7fd5]
 Cho cấp  số cộng $-2$; $x$; $6$; $y$. Khi đó
  \choiceTF
  {\True  $x=2$}
   {$y=8$}
   {$\; P=y-x=6$}
  {\True $P=x^2+y^2=104$} 
  \loigiai{
  \begin{itemchoice}
   \itemch Theo tính chất của cấp số cộng, ta có  $x=\dfrac{-2+6}{2}=2$ và $6=\dfrac{x+y}{2}$.
   \itemch  Vì $x=2$ nên $6=\dfrac{2+y}{2} \Rightarrow y=10$.
   \itemch $\; P=y-x=8$. 
    \itemch $\; P=x^2+y^2=2^2+10^2=104$.
  \end{itemchoice}
 } 
\end{ex}

% nguồn: Ok_dungsai_lop11_Lan1
\begin{ex}%[1D2H2-5]%[ma:3f34304c]
 Cho cấp số cộng $(u_n)$ có $u_2=-2$, công sai $d=2$. Các mệnh đề sau đúng  hay sai?
  \choiceTF
  {Số hạng đầu $u_1$  bằng $4$}
  {\True  Số hạng $u_7$ bằng $8$}
   {\True Số hạng tổng quát $u_n=2n-6$}
  {Tổng $20$ số hạng đầu của cấp  số cộng bằng $400$}
 \loigiai{ 
  \begin{itemchoice}
   \itemch Ta có $u_2-u_1=d\Rightarrow u_1=u_2-d=-2-2=-4$. 
    \itemch Ta có $u_7=u_1+6d=-4+6\cdot2=8$.
   \itemch Số hạng tổng quát $u_n=u_1+(n-1)d=-4+(n-1)\cdot2=2n-6$. 
    \itemch Ta có $S_{n}=nu_1+\dfrac{n(n-1)d}{2}$. Suy ra $S_{20}=20\cdot(-4)+\dfrac{20(20-1)\cdot2}{2}=300$.
  \end{itemchoice}
 } 
\end{ex}

% nguồn: Ok_dungsai_lop11_Lan1
\begin{ex}%[1D2H2-5]%[ma:3cdc074f]
 Cho cấp số cộng $(u_n)$ có $u_1=1$, $u_2=3$. Các mệnh đề sau  đúng hay sai?
  \choiceTF
  {\True Công sai của  cấp số cộng bằng $2$}
  {\True Số hạng  thứ sáu bằng $11$}
   {Số hạng tổng quát là $u_n=2n+1$}
  {\True Tổng $20$  số hạng đầu bằng $400$}
 \loigiai{ 
  \begin{itemchoice}
   \itemch Ta có  $d=u_2-u_1=3-1=2$.
    \itemch Ta có $u_6=u_1+5d=1+5\cdot2=11$
   \itemch Số  hạng tổng quát $u_n=u_1+(n-1)d=1+(n-1)\cdot2=2n-1$.
    \itemch Ta có $S_{n}=nu_1+\dfrac{n(n-1)d}{2}$. Suy ra $S_{20}=20\cdot1+\dfrac{20(20-1)\cdot2}{2}=400$.
  \end{itemchoice}
 } 
\end{ex}

% nguồn: Ok_dungsai_lop11_Lan1
\begin{ex}%[1D2H2-5]%[ma:32bbff0c]
 Cho cấp số cộng $(u_n)$ có $u_2=-1$, $u_3=2$. Các mệnh đề sau đúng hay  sai?
  \choiceTF
  {\True Công sai của cấp  số cộng bằng $3$}
  {Số hạng thứ năm bằng  $11$}
  {\True Số hạng tổng quát là  $u_n=3n-7$}
   {Tổng $25$ số hạng đầu bằng $900$}
  \loigiai{
  \begin{itemchoice}
   \itemch Ta có  $d=u_3-u_2=2-(-1)=3$.
   \itemch Ta có  $u_5=u_2+3d=-1+3\cdot3=8$.
   \itemch Ta có $u_1=u_2-d=-1-3=-4$.\\ 
   Số hạng tổng quát  $u_n=u_1+(n-1)d=-4+(n-1)\cdot3=3n-7$.
   \itemch Ta có $S_{n}=nu_1+\dfrac{n(n-1)d}{2}$.  Suy ra $S_{25}=25\cdot(-4)+\dfrac{25(25-1)\cdot3}{2}=800$.
  \end{itemchoice}
  }
\end{ex}

% nguồn: Ok_dungsai_lop11_Lan1
\begin{ex}%[1D2H2-5]%[ma:b55686fb]
 Cho cấp số cộng  $\left(u_n\right)$, biết $\heva{&u_1+u_2=2\\ &u_5-u_2=24}$. Các mệnh đề sau đúng hay sai?
  \choiceTF
  {Cấp số cộng có $u_1=-2$} 
  {\True  Cấp số cộng có công sai $d=8$}
  {Tổng của $100$ số hạng đầu của  $\left(u_n\right)$ bằng $39000$}
  {\True Số $ 2021$ là số  hạng thứ $254$ của $\left(u_n\right)$}
  \loigiai{
  \begin{itemchoice}
   \itemch Ta có 
   \begin{eqnarray*}
    \heva{&u_1+u_2=2\\ &u_5-u_2=24} &\Leftrightarrow& \heva{&2u_1+d=2\\ &3d=24} \Leftrightarrow  \heva{&u_1=-3\\ &d=8.}
   \end{eqnarray*}
   \itemch Ta có  $S_{n}=nu_1+\dfrac{n(n-1)d}{2}$. Suy ra $S_{100}=100\cdot(-3)+\dfrac{100(100-1)\cdot8}{2}=39300$.
   \itemch Giả sử $2021$ là số hạng thứ $n$ của cấp  số cộng trên. Khi đó
   \begin{eqnarray*}
    u_n=2021 &\Leftrightarrow&  u_1+(n-1)d=2021 \\
    &\Leftrightarrow& -3 + 8(n-1)  = 2021 \Leftrightarrow n= 254.
   \end{eqnarray*}
   Vậy $2021$ là số hạng thứ $254$ của cấp số cộng  trên.
  \end{itemchoice}
 } 
\end{ex}

% nguồn: Ok_dungsai_lop11_Lan1
\begin{ex}%[1D2H2-5]%[ma:c104b386]
 Cho cấp  số cộng $\left(u_n\right)$, biết $\heva{&u_1+3u_5=28\\&S_6-u_3=29}$. Các mệnh đề sau đúng hay sai?
 \choiceTF 
  {\True  Cấp số cộng có $u_1=-2$}
  {\True Cấp số cộng  có công sai $d=3$}
   {Số hạng tổng quát của $\left(u_n\right)$ là $u_n=3n+1$}
   {\True Tổng của $20$ số hạng đầu tiên của cấp số cộng bằng $550$}
  \loigiai{
  \begin{itemchoice}
    \itemch Ta có
   $\heva{&u_1+3u_5=28\\&S_6-u_3=29} \Leftrightarrow  \heva{&u_1+3(u_1+4d)=28\\&(6u_1+15d)-(u_1+2d)=29}\Leftrightarrow \heva{&u_1+3d=7\\&5u_1+13d=29} \Leftrightarrow \heva{&u_1=-2\\&d=3.}$
   \itemch  Ta có $u_n=u_1+(n-1)d=-2+(n-1)\cdot3=3n-5$.
   \itemch  Ta có $S_{20}=\dfrac{20}{2}(2u_1+19d)=550$.
  \end{itemchoice}
 } 
\end{ex}

% nguồn: Ok_dungsai_lop11_Lan1
\begin{ex}%[1D2H2-5]%[ma:9f9b13b7]
 Cho cấp  số cộng $(u_n)$ biết $u_1+2u_4=-21$, $S_{10}=-195$. Các mệnh đề sau đúng hay sai?
 \choiceTF 
  {\True Cấp số cộng có số hạng đầu bằng $3$} 
  {Cấp số cộng có công sai  $d=5$}
   {Số hạng $u_{11}=-50$}
  {Tổng $25$  số hạng đầu của cấp số cộng bằng $-1500$}
 \loigiai{ 
  \begin{itemchoice}
   \itemch Ta có $u_1+2u_4=-21\Leftrightarrow  u_1+2\left(u_1+3d\right)=-21\Leftrightarrow 3u_1+6d=-21.\hfill (1)$\\
   Mặt khác $S_{10}=-195\Leftrightarrow \dfrac{10}{2}\cdot \left(2u_1+9d\right)=-195\Leftrightarrow 2u_1+9d=-39.\hfill (2)$\\ 
   Từ $(1)$ và $(2)$ suy ra $u_1=3$  và $d=-5.$
   \itemch  Do đó $u_{11}=u_1+10d= 1+10\cdot \dfrac(-5)=-49$.
   \itemch Ta có $S_{n}=nu_1+\dfrac{n(n-1)d}{2}$. Suy ra  $S_{25}=25\cdot1+\dfrac{25(25-1)\cdot(-5)}{2}=-1475$.
  \end{itemchoice}
 } 
\end{ex}

% nguồn: Ok_dungsai_lop11_Lan1
\begin{ex}%[1D2H2-5]%[ma:04f70b29]
 Cho cấp số cộng $(u_n)$  có $u_5=15$, công sai $d=3$. Các mệnh đề sau đúng hay sai?
 \choiceTF 
  {\True Số hạng đầu $u_1$  bằng $3$}
   {\True Số hạng $u_6$ bằng $18$}
  {\True Số hạng tổng quát $u_n=3n$} 
  {Tổng $30$ số hạng đầu của  cấp số cộng bằng $1180$}
  \loigiai{
  \begin{itemchoice}
   \itemch Ta  có $u_5=u_1+4d\Leftrightarrow 15=u_1+4\cdot 3\Leftrightarrow u_1=3$.
   \itemch Ta có  $u_6=u_5+d=15+3=18$.
   \itemch Số hạng tổng quát  $u_n=u_1+(n-1)d=3+(n-1)\cdot3=3n$.
   \itemch Ta  có $S_{n}=nu_1+\dfrac{n(n-1)d}{2}$. Suy ra $S_{30}=30\cdot(-3)+\dfrac{30(30-1)\cdot3}{2}=1215$.
  \end{itemchoice}
  }
\end{ex}

% nguồn: Ok_dungsai_lop11_Lan1
\begin{ex}%[1D2H2-5]%[ma:fb011dc2]
 Cho cấp số cộng $(u_n)$ biết $ \heva{& u_3 + u_{10} = -31\\& 2u_{4} - u_9 = 7} $. Các mệnh đề sau đúng  hay sai?
 \choiceTF 
  {Cấp số cộng có  $u_1=-1$, công sai $d=3$}
  {\True Số hạng  $u_3$ bằng $-5$}
   {\True Số hạng tổng quát $u_n=-3n+4$}
  {\True Tổng  $10$ số hạng đầu của cấp số cộng bằng $-125$}
 \loigiai{ 
  \begin{itemchoice}
    \itemch Ta có
   \begin{align*}
    \heva{& u_3 + u_{10} = -31\\& 2u_{4} - u_9 = 7}  \Leftrightarrow \heva{& u_1 + 2d + u_{1} + 9d = -31\\& 2\left( u_{1} + 3d \right) - \left( u_1 + 8d \right) = 7} \Leftrightarrow \heva{& 2u_1 + 11d = -31\\& u_{1} - 2d = 7} \Leftrightarrow \heva{& u_1 = 1\\& d = -3.}
   \end{align*}
   \itemch Ta có  $u_3=u_1+2d=1+2\cdot (-3)=-5$.
   \itemch Số hạng tổng quát  $u_n=u_1+(n-1)d=1+(n-1)\cdot(-3)=-3n+4$.
   \itemch Ta có $S_{n}=nu_1+\dfrac{n(n-1)d}{2}$. Suy ra $S_{10}=10\cdot1+\dfrac{10(10-1)\cdot(-3)}{2}=-125$. 
  \end{itemchoice}
  }
\end{ex}

% nguồn: Ok_dungsai_lop11_Lan2
\begin{ex}%[1D2H2-5]%[ma:f4eb922b]
	Cho cấp số cộng $-2;x;6;y$. Xét tính đúng sai của các khẳng  định sau.
	 \choiceTF
		 {\True $x = 2$}
		{$y =  8$}
		{$P  = y-x = 6$}
		{\True $P = x^2+y^2 = 104$} 
	 \loigiai{
		\begin{itemchoice}
			 \itemch
			Vì Theo tính chất của cấp số cộng, ta có $x = \dfrac{-2+6}{2} =  2$ và $6 = \dfrac{x+y}{2}$.
			\itemch 
			 Vì $x = 2$ nên $6 = \dfrac{2+y}{2} \Rightarrow y = 10$.
			\itemch 
			 Vậy $P = y-x = 8$.
			\itemch 
			Vì $P = x^2+y^2  = 2^2+10^2 = 104$.
		\end{itemchoice}
	 }
\end{ex}


%%%%%%%%%%%%%%%% TRẢ LỜI NGẮN %%%%%%%%%%%%%%%%

% nguồn: Ok_TraLoiNgan_Lop11_Lan1
\begin{ex}%[1D2H2-5]%[ma:701b066d]
	Cho biết  bốn số $5$; $x$; $15$; $y $ theo thứ tự lập thành một cấp số cộng. Tính giá trị của biểu thức $3x+2y$.
	\shortans{70} 
	\loigiai{ 
		 Theo tính chất của cấp số cộng, ta có $\heva{&{x=\dfrac{5+15}{2}} \\
		&{\dfrac{x+y}{2}=15}}\Leftrightarrow \heva{&{x=10} \\ 
		&{y=20}.}$\\ 
		Vậy  $3x+2y=70$.
	 }
\end{ex}

% nguồn: Ok_TraLoiNgan_Lop11_Lan2
\begin{ex}%[1D2H2-5]%[ma:03787fed]
	Câu 13[Dự  án 11 HVA 2024-2025]
	Người ta viết thêm $999$ số thực vào giữa số  $1$ và số $2018$ để được cấp số cộng có $1001$ số hạng. Tìm số hạng thứ $501$.
	 \shortans{1010}
	\loigiai{ 
		Áp dụng công thức cấp số cộng ta  có ${u_n}={u_1}+( n-1)d\Rightarrow {u_{1001}}={u_1}+( 1001-1)d \Rightarrow d=\dfrac{2017}{1000}$.\\
	Vậy số hạng thứ $501$  là ${u_{501}}={u_1}+( 501-1)d=\dfrac{2019}{2}$.}
\end{ex}

% nguồn: Ok_TraLoiNgan_Lop11_Lan2
\begin{ex}%[1D2H2-5]%[ma:b2aa07e8]
	Câu 9[Dự  án 11 HVA 2024-2025]
	Một tam giác vuông có chu vi bằng $9$ và độ dài các cạnh lập thành một cấp số cộng. Diện tích của tam giác  vuông bằng(làm tròn đến hàng phần mười).
	\shortans{1,1} 
	\loigiai{ 
		Gọi $d$ là công sai của cấp số cộng. Độ dài các cạnh của tam giác vuông là $a-d;a;a+d$ với $0 <  d < a$.\\
		Vì tam giác có chu vi bằng $9$  nên $ a-d+a+a+d=9\Leftrightarrow 3a=9\Leftrightarrow a=3$.\\
		Vì tam giác vuông nên  theo Pytago ta có ${{( a-d)}^2}+{a^2}={{( a+d)}^2}\Leftrightarrow {{( 3-d)}^2}+{3^2}={{( 3+d)}^2}\Leftrightarrow 4d={3^2}\Leftrightarrow d=\dfrac{9}{4}$.\\
		Vậy độ dài các cạnh của tam  giác là $\frac{3}{4}$, $3$, $\dfrac{21}{4}$.\\
	Diện tích tam  giác ${S_{\vartriangle }}=\dfrac{1}{2}\cdot \dfrac{3}{4}\cdot 3=\dfrac{9}{8}$.}
\end{ex}

% nguồn: Ok_TraLoiNgan_Lop11_Lan2
\begin{ex}%[1D2H2-5]%[ma:4a456143]
	Câu 14[Dự án 11  HVA 2024-2025]
	Bốn góc của một tứ giác lập thành cấp số cộng. Biết hiệu  giữa góc lớn nhất và góc nhỏ nhất là $6^\circ $. Số đo góc lớn nhất trong các góc đó là (đơn vị độ)
	\shortans{93} 
	\loigiai{ 
		Giả sử $A,B,C,D$ là  $4$ góc của một tứ giác theo thứ tự đó lần lượt lập thành một cấp số cộng và $A<D$.\\
		Ta có $A+B+C+D={360^0}$.\\ 
		Vì  $A,B,C,D$ lập thành cấp số cộng giả sử có công sai $ d>0$, khi đó $B=A+d;C=A+2d;D=A+3d$.\\
		Thay vào ta được $A+( A+d)+( A+2d)+( A+3d)=360^\circ $ $\Leftrightarrow $ $4A+6d=360^\circ  $ $\Leftrightarrow $ $2A+3d=180^\circ $.\\
		Hiệu giữa góc lớn nhất và góc nhỏ nhất là $6^\circ $ nên $D-A=6^\circ $ $\Leftrightarrow $ $3d=6^\circ  $ $\Leftrightarrow $ $ d=2^\circ $.\\
	 Suy ra $A=87^\circ, B=89^\circ, C=91^\circ, D=93^\circ$.}
\end{ex}

% nguồn: Ok_TraLoiNgan_Lop11_Lan2
\begin{ex}%[1D2H2-5]%[ma:db500a19]
	Cho biết bốn số $5;x;15;y$ theo thứ tự lập thành một  cấp số cộng. Tính giá trị của biểu thức $3x+2y$.
	 \shortans{70}
	\loigiai{ 
		Theo tính chất của  cấp số cộng, ta có $\heva{&x =\dfrac{5+15}{2}\\&\dfrac{x+y}{2} = 15} \Leftrightarrow \heva{&x=10\\&y=20.}$\\
		 Vậy $3x+2y=70$.
	 }
\end{ex}

% nguồn: Ok_TraLoiNgan_Lop11_Lan1
\begin{ex}%[1D2V2-5]%[ma:aacbf97b]
	Tìm tất cả các giá trị của tham số  $m$ để phương trình sau có ba nghiệm phân biệt lập thành một cấp số cộng $x^3-3mx^2+2m(m-4)x+9m^2-m=0$.
	 \shortans{1}
	\loigiai{ 
		Phương  trình $x^3-3mx^2+2m(m-4)x+9m^2-m=0$.\hfill{(*)}\\
		Giả sử phương trình đã cho có ba nghiệm phân biệt $x_1$, $x_2$, $x_3$ lập  thành một cấp số cộng.\\
		Theo định lý  Vi-ét đối với phương trình bậc ba, ta có $x_1+x_2+x_3=3m$.\hfill{(1)}\\
		Vì $x_1, x_2, x_3$ lập thành  một cấp số cộng nên $x_1+x_3=2x_2$.\hfill{(2)}\\
		Thay (2) vào (1) ta được  $3x_2=3m\Leftrightarrow x_2=m$.\\
		Thay $x_2=m$  vào phương trình $(*)$ ta được
		\[m^3-3m\cdot m^2+2m(m-4)\cdot m+9m^2-m=0\Leftrightarrow m^2-m=0\Leftrightarrow \hoac{&{m=0}  \\
		 &{m=1}.
		} 
		\] 
		Với $m=0$, ta có $x^3=0\Leftrightarrow x=0$  (loại vì phương trình có một nghiệm duy nhất).\\
		Với $m=1$, ta  có $x^3-3x^2-6x+8=0\Leftrightarrow x=1;x=-2;x=4$ (thỏa mãn).\\
		Vậy $m=1$ thỏa mãn yêu  cầu bài toán.
	} 
\end{ex}

% nguồn: Ok_TraLoiNgan_Lop11_Lan1
\begin{ex}%[1D2V2-5]%[ma:1ba4ba39]
	Có bao nhiêu giá trị $x$ thuộc đoạn $\left[-\pi;\pi \right]$ để $\sin x$; $\sin 2x$; $\sin 3x$ theo thứ tự lập thành cấp số cộng có  công sai $d \ne 0$?
	 \shortans{2}
	 \loigiai{
		Do $\sin x$; $\sin 2x$; $\sin 3x$ nên theo tính chất của cấp số cộng ta  có\\
		$\sin x+\sin 3x=2\sin 2x \Leftrightarrow 2\sin 2x\left(\cos x-1\right)=0 \Leftrightarrow \hoac{&\sin 2x=0 \\& \cos x=1}\Leftrightarrow  x=\dfrac{k\pi}{2}$. \\
		Vì $x \in \left[-\pi;\pi  \right] \Rightarrow -\pi \leq \dfrac{k\pi}{2} \leq \pi \Leftrightarrow -2 \leq k \leq 2$.\\
		Do $d \ne 0$, nên các giá trị  $k\in \left\lbrace-2, 2, 0 \right\rbrace$ không thỏa mãn.\\
	Vậy có $2$ giá  trị $k$ thỏa mãn là $-1$ và $1$.}
\end{ex}

% nguồn: Ok_TraLoiNgan_Lop11_Lan1
\begin{ex}%[1D2V2-5]%[ma:6d9a53f7]
	Tìm tổng tất cả các giá trị $m$ để phương trình $x^4-2(m+1)x^2+2m+1=0$ có bốn nghiệm phân biệt lập thành  cấp số cộng (Kết quả làm tròn đến hàng phần chục).
	 \shortans{3,6}
	\loigiai{ 
		Đặt $t=x^2$, $ t\ge  0$.\\
		Phương trình trở thành $t^2-2(m+1)t+2m+1=0$ \hfill{(2)}\\ 
		Phương trình đề bài có bốn  nghiệm phân biệt khi và chỉ khi phương trình (2) có hai nghiệm dương phân biệt $t_2 > t_1 > 0\Leftrightarrow \heva{&{\Delta ' > 0} \\
		&{P>  0} \\
		&{S> 0}}\Leftrightarrow \heva{&{(m+1)^2-(2m+1) > 0}  \\
		&{2m+1>  0} \\
		&{2(m+1)  > 0}}\Leftrightarrow-\dfrac{1}{2} < m\ne 0$.\\
		Khi đó phương trình (1) có bốn nghiệm  là $-\sqrt{t_2};-\sqrt{t_1};\sqrt{t_1};\sqrt{t_2}$.\\
		Bốn nghiệm này lập thành cấp số cộng khi  $$\heva{&{-\sqrt{t_2}+\sqrt{t_1}=-2\sqrt{t_1}} \\
		 &{-\sqrt{t_1}+\sqrt{t_2}=2\sqrt{t_1}}}\Leftrightarrow \sqrt{t_2}=3\sqrt{t_1} \Leftrightarrow t_2=9t_1.$$
		Theo định lý Vi-ét thì 
		\[\heva{&{t_1+t_2=2(m+1)}  \\
		&{t_1t_2=2m+1}}\Rightarrow  \heva{&{t_1+9t_2=2(m+1)} \\
		&{t_1\cdot  9t_2=2m+1}}\Rightarrow 9m^2-32m-16=0\Leftrightarrow \heva{&m=4\, \text{(thỏa mãn)} \\
		&m=-\dfrac{4}{9} \, \text{(thỏa  mãn)}}
		\] 
		Vậy tổng các giá trị  $m$ cần tìm là $4-\dfrac{4}{9}=\dfrac{32}{9}\approx 3{,}6$.
	 }
\end{ex}


%%%%%%%%%%%%%%%% TỰ LUẬN %%%%%%%%%%%%%%%%

% nguồn: Ok_TuanLuan_Toan11_Lan2
\begin{ex}%[1D2H2-5]%[ma:a3ac1659]
	Biết bốn số $5$; $x$; $15$; $y$ theo thứ tự lập thành cấp số cộng. Giá trị của biểu thức $3x + 2y$ bằng  bao nhiêu?
	\loigiai{Ta có  $x=\dfrac{5+15}{2}=10 \Rightarrow y=20$. Vậy $3 x+2 y=70$.}
\end{ex}

% nguồn: Ok_TuanLuan_Toan11_Lan2
\begin{ex}%[1D2V2-5]%[ma:d8644151]
	Bốn số tạo thành một cấp số cộng có tổng bằng $28$ và tổng  các bình phương của chúng bằng $276$. Tích của bốn số đó là bao nhiêu?
	\loigiai{Gọi  4 số cần tìm là $a-3 r, a-r, a+r, a+3 r$.\\
		Ta có $\heva{&a-3r+a-r+a+r+a+3r=28 \\ &(a-3 r)^{2}+(a-r)^{2}+(a+r)^{2}+(a+3 r)^{2}=276} \Leftrightarrow\heva{&a=7 \\&r^{2}=4}  \Leftrightarrow\heva{&a=7\\ &r=\pm 2.}$\\
	Bốn số cần tìm là $1$, $5$, $9$,  $13$ có tích bằng $585$.}
\end{ex}

% nguồn: Ok_TuanLuan_Toan11_Lan2
\begin{ex}%[1D2V2-5]%[ma:7d8b1cf8]
	Người ta viết thêm $999$ số thực vào giữa  số $1$ và số $2\,018$ để được cấp số cộng có $1\,001$ số hạng. Tìm số hạng thứ $501$.
	\loigiai{Áp dụng công thức cấp số  cộng ta có
		\begin{align*}
			u_{n}=u_{1}+(n-1) d \Rightarrow u_{1\,001}=u_{1}+(1\,001-1) d \Leftrightarrow  2\,018=1+(1\,001-1) d \Rightarrow d=\dfrac{2\,017}{1\,000}.
		\end{align*}
	Vậy số hạng thứ $501$ là  $u_{501}=u_{1}+(501-1) d=\dfrac{2019}{2}$.}
\end{ex}
